\chapter{The transport equation near the separatrix}
\label{chap:transport}

The eigenvalue equation $H(\omega+\psi)=\lambda\omega$ is a first order
transport equation for $\omega$ along the orbits of the dipole, forced by
$\psi$. Near the separatrix the orbits have period $S(h)\to\infty$, and the
equation can be solved explicitly orbit by orbit. The manuscript uses this in
three ways:
\begin{enumerate}
  \item to identify the limit of boundary terms on the separatrix (the
        \emph{canonical transport profile});
  \item to prove weighted estimates for $\omega$ near the saddles (the
        \emph{transport gains});
  \item to prove $H^1$ regularity of $\omega$ across the separatrix.
\end{enumerate}
Each is applied to several different pairs (the eigenfunction itself, the
odd--odd quotient, and a jet-subtracted remainder in the mixed sector). The
manuscript handles this by reusing earlier \emph{proofs}. Here we state the
results for a general transport pair.

\begin{definition}[Transport pair]
\label{def:transport-pair}
\uses{def:eigen-problem, def:lift-operators, def:dipole-flow, def:weak-derivative-local, def:sobolev-local}
\lean{NoEmbeddedEigenvalues.IsOrbitUnion,
  NoEmbeddedEigenvalues.HasWeakL0On,
  NoEmbeddedEigenvalues.IsTransportPair,
  NoEmbeddedEigenvalues.transportUnknown}
\leanok
An \emph{open union of orbits} is an open set $U\subset\R^2$ with
$\varphi_t(U)\subset U$ for all $t$. On such $U$, $L_0g=h$ \emph{weakly} means that
$g,h$ are locally integrable on $U$ and $\int gL_0\phi=-\int h\phi$ for every test
function $\phi$ on $U$. Let $\lambda\in\R\setminus\{0\}$. A \emph{transport pair} on
$U$ is $(v,f)$ with $v\in L^2(U)$, $f\in H^2(U)$ and $H(v+f)=\lambda v$ weakly on $U$,
that is $L_0(v+f)=i\lambda v$, where $H=-iL_0$ as in
Definition~\ref{def:eigen-problem}. Its \emph{transport unknown} is
$z=-(v+f)/\lambda$.
\end{definition}

\begin{remark}
The weak form has no zeroth order term because $Y$ is divergence free, so the
formal adjoint of $L_0$ is $-L_0$. The examples are $(\omega,\psi)$ for an
eigenfunction (Lemma~\ref{lem:transport-pair-eigen}), $(\nu,\phi)$ for the odd--odd
quotient (Proposition~\ref{prop:dd-quotient-domain}), and $(v,f)$ after subtracting
an explicit jet in the mixed sector (Proposition~\ref{prop:mixed-quotient-domain}).

Transport pairs live on the lift $\R^2$, whereas the manuscript works on the
torus. The lift of a periodic function is not in $L^2$ of an unbounded set, so in
the applications $U$ is bounded. Such $U$ exist: every orbit away from the saddles
is a closed loop in one cell, and every separatrix orbit is a single bounded arc, so
for instance the closed positive cell together with thin annuli of orbits inside
the neighbouring cells is a bounded open union of orbits.
\end{remark}

\begin{lemma}[The transport unknown]
\label{lem:transport-unknown}
\uses{def:transport-pair}
\lean{NoEmbeddedEigenvalues.HasWeakL0On.const_mul,
  NoEmbeddedEigenvalues.IsTransportPair.hasWeakL0On_transportUnknown,
  NoEmbeddedEigenvalues.IsTransportPair.hasWeakL0On_transportUnknown_sub}
\leanok
For a transport pair, the transport unknown satisfies $Hz=-v$ and $(H-\lambda)z=f$
weakly on $U$.
\end{lemma}

\begin{proof}
\leanok
The weak relation is linear, and $z=-(v+f)/\lambda$ gives
$L_0z=-i v$ and $i(\lambda z+f)=-iv$.
\end{proof}

\begin{lemma}[Eigenfunctions give transport pairs]
\label{lem:transport-pair-eigen}
\uses{def:transport-pair, def:eigen-problem, def:torus-of-lift}
Let $\omega$ be an eigenfunction with parameter $\lambda$, $\psi=\Delta^{-1}\omega$,
and $U$ a bounded open union of orbits. Then the lifts of $\omega$ and $\psi$ form a
transport pair on $U$.
\end{lemma}

\begin{proof}
\uses{lem:elliptic-torus, lem:lift-torus, lem:local-to-global}
$\psi\in H^2(\T^2)$, so its lift is in $H^2(U)$ on the bounded set $U$, and the lift
of $\omega$ is in $L^2(U)$. The torus equation is tested against characters; by
density it holds against smooth periodic functions, and a test function on $U$
periodizes to one, which turns the torus equation into the weak equation on $U$.
\end{proof}

\begin{lemma}[Orbitwise transport and resonant vanishing]
\label{lem:resonant-vanishing}
\uses{def:transport-pair, def:period, lem:period-orbit, def:h1-interval, def:h1-periodic}
Let $(v,f)$ be a transport pair on a compact regular orbit annulus. On almost
every orbit, $z$ restricted to the orbit is in $H^1$ of flow time and
satisfies $(-i\partial_t-\lambda)z=f$ there. In angular Fourier modes with
respect to the orbital frequency $b(h)$,
\[
  (kb(h)-\lambda)\,v_k(h)=-kb(h)f_k(h).
\]
If $kb(h_k)=\lambda$ at a level $h_k$ in the annulus, then $f_k(h_k)=0$.
\end{lemma}

\begin{proof}
\uses{lem:coarea, lem:period, lem:sobolev-embedding, lem:local-to-global}
On orbits, $L_0=\partial_t$. The coefficient $f_k$ is continuous in $h$, and
$v_k$ is locally square integrable in $h$. Since $S'<0$ the denominator
$kb(h)-\lambda$ has a simple zero at $h_k$, so $f_k(h_k)\ne0$ would make $v_k$
behave like $1/(h-h_k)$.
\end{proof}

\section{The canonical transport profile on a separatrix arc}

Throughout this section $(v,f)$ is a transport pair on a neighbourhood of the
closed positive cell that is invariant under $R$ with character
$\kappa\in\{1,-1\}$: $R^*f=\kappa f$, $R^*v=\kappa v$. Let
$g(s)=f(\gamma(s))$ be the trace of $f$ on the arc of
Definition~\ref{def:separatrix-arc}, with endpoint limits $c_\pm$. Then
$c_+=\kappa c_-$.

\begin{lemma}[Melnikov condition]
\label{lem:melnikov}
\uses{def:transport-pair, def:separatrix-arc, lem:resonant-vanishing}
$\int_\R g'(s)e^{-i\lambda s}\,ds=0$ (manuscript \texttt{(ee:melnikov)}).
\end{lemma}

\begin{proof}
\uses{lem:resonant-vanishing, lem:period, lem:orbit-comparison, lem:traces}
There are infinitely many resonant levels $h_k\to0$ with $kb(h_k)=\lambda$ and
$e^{i\pi k/2}=\kappa$. By the character, the four quarter-orbit integrals of
$f_k$ coincide, so
$\int_{-L}^{L}f(\gamma_{h_k}(t))e^{-i\lambda t}\,dt=0$. Integration by parts
(the endpoint terms cancel because $e^{2i\lambda L(h_k)}=\kappa$) gives the
same for the derivative. Compare with $g'$ using
$\|g_h-g\|_{L^\infty(-L,L)}\le C_\alpha h^{\alpha/2}$ and
$\|g_h-g\|_{L^1(-L,L)}\le C_\alpha h^{\alpha/2}|\log h|$
(manuscript \texttt{(ee:forcing-comparison)}). Finally $g'\in L^1(ds)$ since
the trace of an $H^2$ function on a straight arc is $H^{3/2}$.
\end{proof}

\begin{definition}[Canonical transport profile]
\label{def:canonical-profile}
\lean{NoEmbeddedEigenvalues.profileIntegrand,
  NoEmbeddedEigenvalues.canonicalProfile,
  NoEmbeddedEigenvalues.canonicalUnknown}
\leanok
Let $\lambda\in\R$ and $g\colon\R\to\C$ with $g'\in L^1(\R)$; in the applications $g$ is
the trace on a separatrix arc (Definition~\ref{def:separatrix-arc}). The
\emph{canonical profile} of $g$ is
\[
  \Omega(s)=e^{i\lambda s}\int_s^\infty g'(t)e^{-i\lambda t}\,dt,
\]
and the \emph{canonical transport unknown} is $z=-(\Omega+g)/\lambda$. Write
$P=-i\partial_s$. The \emph{Melnikov condition} is
$\int_\R g'(t)e^{-i\lambda t}\,dt=0$.
\end{definition}

\begin{remark}[The profile is defined by a formula]
The manuscript defines $\Omega$ as the solution of
$(\partial_s-i\lambda)\Omega=-g'$ with $\Omega(s)\to0$ as $s\to\pm\infty$
(manuscript \texttt{(ee:canonical-ode)}), which presupposes existence. Here
$\Omega$ is the solution that decays at $+\infty$, which always exists; it decays at
$-\infty$ exactly under the Melnikov condition
(Lemma~\ref{lem:canonical-profile-ode}). This makes explicit that the Melnikov
condition (Lemma~\ref{lem:melnikov}) is what the canonical profile needs. In Lean,
$g'$ is the pointwise derivative \texttt{deriv g}, which for locally absolutely
continuous $g$ represents the weak derivative.
\end{remark}

\begin{lemma}[The profile equation]
\label{lem:canonical-profile-ode}
\uses{def:canonical-profile}
\lean{NoEmbeddedEigenvalues.integrable_profileIntegrand,
  NoEmbeddedEigenvalues.ae_hasDerivAt_canonicalProfile,
  NoEmbeddedEigenvalues.ae_hasDerivAt_canonicalUnknown,
  NoEmbeddedEigenvalues.ae_hasDerivAt_canonicalUnknown_sub,
  NoEmbeddedEigenvalues.tendsto_canonicalProfile_atTop,
  NoEmbeddedEigenvalues.tendsto_canonicalProfile_atBot_iff,
  NoEmbeddedEigenvalues.canonicalProfile_eq_Iic}
\leanok
Let $g'\in L^1(\R)$. Then $\Omega'=i\lambda\Omega-g'$ almost everywhere. If moreover
$\lambda\ne0$ and $g$ is differentiable almost everywhere, then $Pz=-\Omega$ and
$(P-\lambda)z=g$ almost everywhere. Further $\Omega(s)\to0$ as $s\to+\infty$, and
$\Omega(s)\to0$ as $s\to-\infty$ if and only if the Melnikov condition holds; in that
case $\Omega(s)=-e^{i\lambda s}\int_{-\infty}^sg'(t)e^{-i\lambda t}\,dt$.
\end{lemma}

\begin{proof}
\leanok
Write $J(s)=\int_s^\infty q$ with $q(t)=g'(t)e^{-i\lambda t}\in L^1$. By Lebesgue
differentiation $J'=-q$ almost everywhere, so
$\Omega'=i\lambda\Omega+e^{i\lambda s}J'=i\lambda\Omega-g'$. Then
$z'=-(\Omega'+g')/\lambda=-i\Omega$, that is $Pz=-\Omega$, and
$(P-\lambda)z=-\Omega+\Omega+g=g$. Since $|\Omega|=|J|$, $J(s)\to0$ as $s\to+\infty$
and $J(s)\to\int_\R q$ as $s\to-\infty$. Finally
$\int_{-\infty}^sq+\int_s^\infty q=\int_\R q$.
\end{proof}

\begin{lemma}[Uniqueness and decay of the canonical profile]
\label{lem:canonical-profile}
\uses{def:canonical-profile, lem:melnikov}
Under the Melnikov condition, $\Omega$ is the unique locally absolutely continuous
solution of $(\partial_s-i\lambda)\Omega=-g'$ with $\Omega(s)\to0$ as
$s\to\pm\infty$. For the arc trace $g$ of an $H^2$ function it satisfies
$|\Omega(s)|\le C_\alpha e^{-\alpha|s|}$ for every $\alpha<1$, and
$\|e^{|s|}\Omega\|_{L^2(ds)}\le C\|e^{|s|}g'\|_{L^2(ds)}$ whenever the right
side is finite.
\end{lemma}

\begin{proof}
\uses{lem:canonical-profile-ode, lem:melnikov, lem:sobolev-embedding, lem:local-to-global}
Existence is Lemma~\ref{lem:canonical-profile-ode}. A difference of two solutions is
$Ce^{i\lambda s}$, which does not decay. Decay follows by integration by parts in
the two formulas of Lemma~\ref{lem:canonical-profile-ode}, using
$g-c_\pm=O(e^{-\alpha|s|})$, which comes from $H^2\subset C^{0,\alpha}$. The weighted
bound is a one-sided convolution estimate with kernel $e^{-t}\mathbf 1_{t>0}$.
\end{proof}

\begin{lemma}[Good windows]
\label{lem:good-windows}
\uses{def:period}
The set of levels $h$ with $|\lambda S(h)-2\pi(j+\tfrac12)|<\pi/3$ for some
integer $j$ (of the sign of $\lambda$) has positive measure in every interval
$(0,\eta)$, and on it $|e^{i\lambda S(h)/4}-\kappa|\ge c>0$ uniformly
(manuscript \texttt{(ee:good-windows)}).
\end{lemma}

\begin{proof}
\uses{lem:period}
$S$ is continuous, strictly monotone and tends to $\infty$ at $0$. On the
windows, $\lambda S/4$ lies within $\pi/12$ of an odd multiple of $\pi/4$.
\end{proof}

\begin{lemma}[Limit of the vorticity pairing]
\label{lem:pairing-limit}
\uses{def:canonical-profile, def:transport-pair, lem:good-windows}
There are levels $h_j\to0$ in the good windows, avoiding any prescribed null
set, such that
\[
  \int_{\{\Psi=2h_j\}}\operatorname{Re}(v\bar f)\,dt
  \longrightarrow\sum_\gamma\int_\R\operatorname{Re}(\Omega_\gamma\bar g_\gamma)\,ds
\]
where the sum runs over the four arcs bounding $\Cell$ (manuscript
\texttt{(ee:pairing-limit)}).
\end{lemma}

\begin{proof}
\uses{lem:canonical-profile, lem:canonical-profile-ode, lem:resonant-vanishing, lem:orbit-comparison, lem:good-windows}
On a good orbit $z_h=-(v_h+g_h)/\lambda$ solves $(P-\lambda)z_h=g_h$ with
$z_h(L)=\kappa z_h(-L)$. The difference $z_h-z$ solves the same equation with
forcing $g_h-g$ and boundary mismatch $\kappa z(-L)-z(L)=O(e^{-\alpha L})$
(using $c_+=\kappa c_-$). Variation of constants and the window lower bound
give $\|v_h-\Omega\|_{L^\infty(-L,L)}\le C_\alpha h^{\alpha/2}|\log h|$. Since
$L=O(|\log h|)$, the pairing error is $O(h^{\alpha/2}|\log h|^2)$, and the
tails of $\Omega\bar g$ vanish by exponential decay.
\end{proof}

\section{Weighted transport gains near the saddles}

Let $d$ denote distance to the saddle set. In this section $(v,f)$ is a
transport pair on a neighbourhood of the separatrix, decomposed into rotation
characters as in Lemma~\ref{lem:rotation-characters}.

\begin{definition}[Quarter-orbit numerator and divisor]
\label{def:numerator}
\uses{def:transport-pair, def:period, lem:period-orbit}
For the quarter orbits of Definition~\ref{def:period} and a character $\kappa$, let
\[
  N(h)=\int_{-L(h)}^{L(h)}f(\gamma_h(t))e^{-i\lambda t}\,dt,\qquad
  D(h)=\kappa-e^{2i\lambda L(h)},\qquad C(h)=N(h)/D(h).
\]
Variation of constants gives $|z_h(-L)|=|C(h)|$ and $z_h(L)=\kappa z_h(-L)$
(manuscript \texttt{(eq:quarter-transport-divisor)}).
\end{definition}

\begin{lemma}[Numerator vanishes at divisor zeros]
\label{lem:numerator-zeros}
\uses{def:numerator, def:h1-interval}
$N\in H^1_{\mathrm{loc}}(0,h_0)$ and $N(h_j)=0$ at every zero $h_j$ of $D$.
The zeros are simple and form a geometric sequence. On the phase cell $I_j$
around $h_j$ one has $h\sim h_j$ and $|D(h)|\ge c_\lambda|h-h_j|/h_j$.
\end{lemma}

\begin{proof}
\uses{lem:period, lem:resonant-vanishing, lem:coarea}
The endpoint traces of $z_h$ are locally square integrable in $h$, so
$|N/D|$ cannot blow up like $1/|h-h_j|$; continuity of $N$ forces $N(h_j)=0$.
Simplicity and spacing follow from $L'(h)\sim-1/(2h)$.
\end{proof}

\begin{lemma}[First transport gain]
\label{lem:transport-first-gain}
\uses{def:transport-pair}
\notready
If $(v,f)$ is a transport pair near the separatrix with $f\in H^2$, then
$v/d\in L^2$ near the saddles. No elliptic relation between $v$ and $f$ is
used.
\end{lemma}

\begin{proof}
\uses{lem:numerator-zeros, lem:orbit-comparison, lem:hardy-1d, lem:sobolev-embedding, lem:local-to-global}
First assume $f$ vanishes at the saddles. Differentiating $N$ (including the
moving endpoints) gives
$|N'(h)|^2\le Ch^{-1}\int_{-L}^{L}|\nabla f(\gamma_h)|^2\,dt+Ch^{\alpha-2}$
(manuscript \texttt{(eq:first-numerator-estimate)}). Hardy's inequality on
each side of $h_j$ in each phase cell gives
$\int_{I_j}|N/D|^2\,dh/h\le C\int_{I_j}\!\int|\nabla f_h|^2\,dt\,dh+Ch_j^\alpha$,
which is summable. Reconstruction of $v_h$ from the endpoint value
gives $|v_h(t)|\le C|N/D|+Ce^{-\alpha|t|}$, and
$\int_{-L}^Ld^{-2}dt\le C/h$ completes the estimate. For nonzero saddle values
subtract a smooth function $a(t)$ with the same tail values (these satisfy
$a_+=\kappa a_-$); the same estimates apply to the difference.
\end{proof}

\begin{lemma}[Second transport gain]
\label{lem:transport-second-gain}
\uses{def:transport-pair}
\notready
Let $0<\delta<1$. If in addition $|f|\le Cd^{2-\delta}$ near the saddles and
$\int_0^{h_0}\!\int_{-L}^{L}|\nabla f(\gamma_h(t))|^2d^{-(4-2\delta)}\,dt\,dh<\infty$,
then $v/d^2\in L^2$ near the saddles.
\end{lemma}

\begin{proof}
\uses{lem:transport-first-gain, lem:numerator-zeros, lem:orbit-comparison, lem:hardy-1d}
Now $|N'|^2\le C\int|\nabla f_h|^2d^{-(4-2\delta)}\,dt+Ch^{-\delta}$, so
$N'\in L^2$, and the phase-cell Hardy inequality with weight $h^{-2}$ gives
$\int_{I_j}|N/D|^2\,dh/h^2\le C\int_{I_j}|N'|^2$. Reconstruction has forcing
tail $Ce^{-(2-\delta)|t|}$, whose squared $d^{-4}$ integral is $O(h^{-\delta})$,
while $\int d^{-4}dt=O(h^{-2})$.
\end{proof}

\begin{lemma}[$H^1$ transport gain]
\label{lem:transport-h1-gain}
\uses{def:transport-pair}
\notready
Suppose in addition that $hN''\in L^2(0,h_0)$, $C/h\in L^2(0,h_0)$, and
$\int_0^{h_0}\!\int_{-L}^L|\nabla f_h|^2d^{-(4-2\delta)}\,dt\,dh<\infty$. Then
$d\,\partial_hv$ and $d^{-1}\partial_tv$ are square integrable in the orbit
coordinates, and hence $v\in H^1$ on each cell's separatrix neighbourhood.
If moreover the midpoint value $A(h)=z_h(0)$ satisfies $A(0)=0$, the trace of
$v$ on the regular edges is zero.
\end{lemma}

\begin{proof}
\uses{lem:numerator-zeros, lem:hardy-1d, lem:orbit-comparison}
Factor $D=(h-h_j)d_j(h)/h_j$ on each phase cell. The second Hardy inequality
gives $C'\in L^2$. The midpoint formula
$A(h)=ie^{2i\lambda L}C+i\int_{-L}^0f_he^{-i\lambda t}dt$ is then in $H^1$,
and $z_h(t)=e^{i\lambda t}A(h)+i\int_0^te^{i\lambda(t-u)}f_h(u)\,du$
controls $\partial_hz$. Finally $\partial_tv=i\lambda v-\partial_tf$, and the
inverse coordinate bounds $|\nabla h|\le Cd$, $|\nabla t|\le C/d$ convert to
physical derivatives.
\end{proof}

\begin{remark}[Digest note]
Lemmas~\ref{lem:transport-first-gain}--\ref{lem:transport-h1-gain} are
extracted from the proofs of the manuscript's
Propositions~\texttt{prop:saddle-gains}, \texttt{prop:separatrix-h1},
\texttt{prop:dd-quotient-domain} and \texttt{prop:mixed-quotient-domain}. The
hypotheses above are a first guess at what those proofs actually use and must
be pinned down before formalization; hence \verb|\notready|. Getting these
statements right is probably the single most useful restructuring for the
paper as well.
\end{remark}
